\section{Correlated Pell--Lucas staircases and opposite-channel incidence}
\label{sec:pell-lucas-correlated-all-order}

We continue the norm-one splice of
Section~\ref{sec:pell-lucas-all-order-staircase}.  For a norm-one Lucas pair
write
\[
 u_n=\frac{\alpha^n-\beta^n}{\alpha-\beta},
 \qquad v_n=\alpha^n+\beta^n,
 \qquad \alpha\beta=1,
\]
and let $\delta=(\alpha-\beta)^2$.  For a positive odd integer
$k=2\theta+1$, define
\begin{equation}\label{eq:companion-polynomial-coefficients}
 d_j(k)=\binom{\theta+j}{2j},
 \qquad
 \Psi_k(W)=\sum_{j=0}^{\theta}d_j(k)W^j.
\end{equation}

\begin{proposition}[Exact companion multiplication polynomial]
\label{prop:exact-companion-multiplication-polynomial}
For every $n\ge1$,
\begin{equation}\label{eq:exact-companion-multiplication-polynomial}
 v_{nk}=v_n\Psi_k(\delta u_n^2).
\end{equation}
Thus the notation $v_{nk}/v_n$ below denotes an exact integer identity and
does not assume division in a residue ring.
\end{proposition}

\begin{proof}
In a splitting ring put $x=\alpha^n$ and
$W=(x-x^{-1})^2=\delta u_n^2$.  No quotient in the splitting ring is
needed.  Put $S_m=x^{2m+1}+x^{-(2m+1)}$.  The exact Laurent-polynomial
recurrence is $S_{m+1}=(W+2)S_m-S_{m-1}$, with
$S_0=x+x^{-1}$ and $S_1=(x+x^{-1})(1+W)$.  The polynomials
$P_m(W)=\sum_{j=0}^m\binom{m+j}{2j}W^j$ have the same recurrence and
initial values, by Pascal's identity.  Induction therefore gives the exact
factorization $S_m=(x+x^{-1})P_m(W)$ without cancelling a possibly nonunit
factor.  Taking $m=\theta$ is
\eqref{eq:exact-companion-multiplication-polynomial}; being symmetric in
$x,x^{-1}$, it descends to the integer identity.
\end{proof}

Let $\Phi_k(W)=\sum c_j(k)W^j$ be the first-sequence multiplication
polynomial from \cite[Proposition~5.1]{ZhangLucas2026}, where
\[
 c_j(k)=\frac{k\prod_{h=1}^j(k^2-(2h-1)^2)}
              {4^j(2j+1)!}.
\]

\begin{theorem}[Coefficientwise companion correlation]
\label{thm:pell-lucas-coefficient-correlation}
For $0\le j\le\theta$,
\begin{equation}\label{eq:pell-lucas-coefficient-correlation}
             (2j+1)c_j(k)=k d_j(k).
\end{equation}
Consequently
\begin{equation}\label{eq:pell-lucas-differential-polynomials}
                 \Phi_k(W)+2W\Phi_k'(W)=k\Psi_k(W).
\end{equation}
\end{theorem}

\begin{proof}
Since $k=2\theta+1$,
\[
 \frac1{4^j}\prod_{h=1}^j(k^2-(2h-1)^2)
 =\frac{(\theta+j)!}{(\theta-j)!}.
\]
Substitution gives
$c_j(k)=\frac{k}{2j+1}\binom{\theta+j}{2j}$, proving
\eqref{eq:pell-lucas-coefficient-correlation}.  Coefficient comparison then
proves \eqref{eq:pell-lucas-differential-polynomials}.
\end{proof}

Specialize to the balancing-Pell pair at an odd prime index $\ell$:
\[
 A=A_\ell,\qquad B=B_\ell,\qquad U=AB=u_\ell,\qquad
 v_\ell=2A_{2\ell},\qquad \delta=32.
\]
Put $\theta=(\ell-1)/2$, $a_j=32^jc_j(\ell)$,
$b_j=32^jd_j(\ell)$ and, for
$0\le s\le(\ell-1)/2$, define
\[
 \mathcal E_s(X)=\sum_{j=s}^{\theta}a_jX^{j-s},\qquad
 \mathcal F_s(X)=\sum_{j=s}^{\theta}b_jX^{j-s},\qquad
 E_s=\mathcal E_s(U^2),\quad F_s=\mathcal F_s(U^2).
\]

\begin{theorem}[Paired all-order support-unit staircase]
\label{thm:pell-lucas-correlated-staircase}
Removing the first $s$ terms from either multiplication polynomial leaves
respectively $U^{2s}E_s$ and $U^{2s}F_s$.  Moreover
\begin{equation}\label{eq:pell-lucas-tail-units-and-correlation}
 \gcd(E_s,U)=\gcd(F_s,U)=1,
 \qquad
 (2s+1)\mathcal E_s(X)+2X\mathcal E_s'(X)
       =\ell\mathcal F_s(X),
\end{equation}
and in particular
\begin{equation}\label{eq:pell-lucas-tail-mod-correlation}
                         (2s+1)E_s\equiv\ell F_s\pmod{U^2}.
\end{equation}
If $p\mid U$ and $e_p=v_p(U)$, both unnormalized tails have exact valuation
$2se_p$ at $p$.
\end{theorem}

\begin{proof}
The tail factorizations are exact by construction.  Every support prime
$p\mid U$ satisfies $p\ge2\ell-1>\ell$.  All factorial arguments in
$d_s(\ell)=\binom{(\ell-1)/2+s}{2s}$ are therefore below $p$, so $d_s$ is a
$p$-adic unit.  The first-sequence support-unit argument from the preceding
section makes $c_s$ a $p$-adic unit, and $p\nmid32$.  Hence the leading terms
$a_s,b_s$, and thus $E_s,F_s$ modulo $U^2$, are units at every support prime.

Apply \eqref{eq:pell-lucas-coefficient-correlation} term by term.  The
coefficient of $X^{j-s}$ on the left of the differential identity is
$(2s+1+2(j-s))a_j=(2j+1)a_j=\ell b_j$.  This proves
\eqref{eq:pell-lucas-tail-units-and-correlation} and its reduction
\eqref{eq:pell-lucas-tail-mod-correlation} after setting $X=U^2$.  The exact
valuations follow from the unit leading tails.
\end{proof}

Set $T_s=v_\ell F_s$ and $Z_0=v_\ell/2=A_{2\ell}$.

\begin{theorem}[Coherent every-order channel splitters]
\label{thm:pell-lucas-every-order-splitters}
For every $s$, the congruence
\begin{equation}\label{eq:pell-lucas-every-order-reconstruction}
 2(2s+1)E_s Z_s\equiv\ell T_s\pmod{U^2}
\end{equation}
has the unique solution
$Z_s\equiv Z_0\pmod{U^2}$.  Consequently
\begin{equation}\label{eq:pell-lucas-every-order-signs}
 Z_s\equiv1\pmod{A^2},\qquad
 Z_s\equiv-1\pmod{B^2},\qquad Z_s^2\equiv1\pmod{U^2}.
\end{equation}
For two levels $s,t$ one has
\begin{equation}\label{eq:pell-lucas-cross-order-determinant}
 (2t+1)E_tT_s\equiv(2s+1)E_sT_t\pmod{U^2}.
\end{equation}
If $q^3\mid A$ and $r^3\mid B$, every $Z_s$ satisfies
\begin{equation}\label{eq:pell-lucas-opposite-depth-six}
                 Z_s\equiv1\pmod{q^6},\qquad
                 Z_s\equiv-1\pmod{r^6}.
\end{equation}
\end{theorem}

\begin{proof}
Multiplying \eqref{eq:pell-lucas-tail-mod-correlation} by $v_\ell$ gives
\[
 \ell T_s\equiv(2s+1)v_\ell E_s
              =2(2s+1)E_sZ_0\pmod{U^2}.
\]
Every prime divisor of $U$ exceeds $\ell$, so
$2(2s+1)E_s$ is invertible modulo $U^2$.  This proves uniqueness and
identifies $Z_s$.  The Pell identities
$A_{2\ell}=2A^2+1=4B^2-1$ give the two signs; coprimality of $A,B$ combines
them modulo $U^2$.  Cross-multiplication at levels $s,t$ and cancellation of
$\ell$, which is coprime to $U$, proves
\eqref{eq:pell-lucas-cross-order-determinant}.  Finally $q^3\mid A$ and
$r^3\mid B$ imply $q^6\mid A^2$ and $r^6\mid B^2$.
\end{proof}

The splitters therefore form one coherent projective line; higher orders are
not independent tests.  The character information can nevertheless be
resolved vertex by vertex.  In a hypothetical squarefull packet write
\[
 A=D_A^3X^2,\qquad B=D_B^3Y^2,
\]
with $D_A,D_B$ squarefree, and let $O_A,O_B$ be their prime supports.

\begin{theorem}[Vertexwise character incidence]
\label{thm:pell-lucas-vertexwise-incidence}
For $r\in O_B$ and $q\in O_A$ one has
\begin{equation}\label{eq:pell-lucas-row-column-laws}
 \prod_{q\in O_A}\left(\frac qr\right)=\left(\frac2r\right),
 \qquad
 \prod_{r\in O_B}\left(\frac qr\right)=\left(\frac Bq\right).
\end{equation}
Both total products equal $(2/\ell)$.  If $q\equiv1\pmod8$, the second
endpoint is the quartic-two sign $2^{(q-1)/4}\pmod q$.
For $\ell\equiv3$ or $5\pmod8$, there exist $q\in O_A,r\in O_B$ such that
\begin{equation}\label{eq:pell-lucas-forced-opposite-pair}
 v_q(A)\ge3,\quad v_r(B)\ge3,\quad
 \left(\frac qr\right)=-1,\quad 2\ell\mid q+r,
\end{equation}
and the every-order signs \eqref{eq:pell-lucas-opposite-depth-six} hold at
this same pair.
\end{theorem}

\begin{proof}
For $r\mid B$, the Pell equation gives $A^2\equiv-1\pmod r$.  Every
$B$-channel prime is $1\pmod4$, so Euler's criterion yields
$(A/r)=(2/r)$.  Removing the square core from
$A=D_A^3X^2$ proves the first product.  Removing the square core from $B$ and
using quadratic reciprocity, whose sign is trivial because $r\equiv1\pmod4$,
proves the second.  Multiplying all rows or all columns gives the same edge
product, and the inherited kernel ledger evaluates it as $(2/\ell)$.

If $q\equiv1\pmod8$, the congruence $2B^2\equiv1\pmod q$ implies
$(B/q)=2^{(q-1)/4}$ because this quarter-power is a sign.  Finally, in the
two stated index classes the inherited table gives $D_B\equiv5\pmod8$.
Thus some $r\in O_B$ is $5\pmod8$ and its row product is negative; one edge
$(q/r)$ is negative.  The rank congruences are
$q\equiv1\pmod{2\ell}$ and $r\equiv-1\pmod{2\ell}$, proving the last
divisibility.  Squarefull odd exponents are at least three, and
Theorem~\ref{thm:pell-lucas-every-order-splitters} supplies the depth-six
signs.
\end{proof}

This necessary packet is highly constrained but has not been contradicted.
An independently replayed computation proves a separate finite theorem:
for every one of the $57$ odd prime indices $3\le\ell\le271$, at least one
of $A_\ell,B_\ell$ has a certified prime divisor of exponent exactly one.
Hence $A_\ell B_\ell$ is not squarefull throughout that range.  The verifier
checks each divisor, its primality, and the coordinate modulo its square; the
large witness at $\ell=59$ uses a complete Pocklington certificate.  It also
checks $138{,}675$ coefficient pairs, $228$ recurrence evaluations, and the
incidence graphs at all $13$ prime indices through $43$.

A bounded search of $527{,}352$ candidate primes up to $2\cdot10^6$ finds
$13^2\parallel B_7$ but no depth-three opposite pair.  This is only a finite
no-hit and retires no assertion.  At index $11$, the row at $r=5741$ contains
both $(23/r)=-1$ and $(353/r)=+1$, a full-premise counterexample to the
stronger claim that every edge of a negative row is negative.  It is not a
squarefull packet and therefore does not refute the active exclusion route.

Any surviving packet must simultaneously satisfy the earlier third-order
congruence modulo $8\ell^3$, both support-unit staircases, every cross-order
determinant, the depth-six signs, and the vertexwise incidence laws.  The
missing theorem must couple the prime-factor quotients in the third-order
ledger to this all-order projective line.  Its current absence is a difficult
open gate, not evidence against the route.

The Lean companion checks the rational coefficient identity, exact weighted
list identity, normalized-tail congruence, every-order recovery, cross-order
determinant, sixth-power propagation, negative-row extraction, and the
quartic-two power algebra.  The Lucas multiplication theorem, rank theorem,
finite factor certificates, and the open uniform exclusion are kept at their
stated mathematical or computational boundaries and are not added as axioms.
